\documentclass[10pt,a4paper]{amsart}
\usepackage[margin=19mm]{geometry}
\usepackage{amsmath,amssymb,mathtools}
\usepackage[scr=rsfs,cal=euler]{mathalfa}
\usepackage{xfrac}
\usepackage[unicode,hidelinks,bookmarks=false]{hyperref}
\newcommand{\ie}{i.e.\ }
\newcommand{\cf}{cf.\ }
\newcommand{\resp}{respectively\ }
\newcommand{\Subsection}{\subsection}
\providecommand{\nobreakdash}{\nobreak\hbox{-}}
\setlength{\emergencystretch}{2em}
\multlinegap0pt
\usepackage{latexsym}
\usepackage{float}
\usepackage{tikz}
\usepackage{circuitikz}
\usepackage{mathtools}
\usepackage{enumitem}
\allowdisplaybreaks
\newtheorem{theorem}{Theorem}[section]      
\newtheorem{lemma}[theorem]{Lemma}               
\newtheorem{corollary}[theorem]{Corollary}
\newtheorem{definition}[theorem]{Definition}
\newtheorem{assumption}[theorem]{Assumption}
\newtheorem{example}[theorem]{Example}
\newtheorem{exercise}[theorem]{Exercise}
\newtheorem{remark}[theorem]{Remark}
\newtheorem{proposition}[theorem]{Proposition}
\newtheorem{claim}[theorem]{Claim}
\newtheorem{conjecture}[theorem]{Conjecture}
\DeclareMathOperator\supp{supp}
\DeclareMathOperator\dist{dist}
\DeclareMathOperator{\sgn}{sgn}
\renewcommand{\a}{a}%admittance
\renewcommand{\Re}{\operatorname{Re}}%real part
\renewcommand{\Im}{\operatorname{Im}}%imaginary part
\newcommand{\Deg}{\operatorname{Deg}}
\newcommand{\cp}{\operatorname{cap}}
\newcommand{\abs}[1]{\left|#1\right|}%imaginary part
\newcommand{\F}{\mathcal{F}}
\renewcommand{\L}{\mathcal{L}}
\newcommand{\D}{\mathcal{D}}
\newcommand{\Q}{\mathcal{Q}}
\newcommand{\cfreq}{s}
\begin{document}
\title[m-sectorial discrete Laplacians and recurrence of complex-we]{$m$-sectorial discrete Laplacians and recurrence of complex-weighted graphs}
\author{Anna Muranova}
\date{}
\maketitle
\noindent\emph{Source and licence.} $m$-sectorial discrete Laplacians and recurrence of complex-weighted graphs. Version arXiv:2606.16554v1 (CC BY 4.0). This material is distributed under the Creative Commons Attribution 4.0 International licence, http://creativecommons.org/licenses/by/4.0/, as recorded in the arXiv OAI licence metadata for this exact version and shown on the abstract page https://arxiv.org/abs/2606.16554v1. Full rights notice: licenses/arxiv-2606.16554-CC-BY-4.0.txt.\par\smallskip
\noindent\textit{Editorial scope.} Counted unit: the original \texttt{theorem} environment \texttt{-} at source lines 285--305 together with its complete proof at lines 307--309; the delivered document keeps source lines 118--314 so that every referenced label and every citation resolves inside it. Native editable source is the paper's own concatenated TeX; the AMS wrapper is editorial formatting only. No publisher class, upstream script or unrelated artwork is redistributed.
\section{Preliminaries}
\label{sect::preliminaries}
\subsection{Complex-weighted graphs and formal Laplacian}

\begin{definition}\label{def::graph}
A {\em graph (with a measure)} is a triple $(X,E,m)$, where $X$ is an at most countable set, 
$$
E
 \subset \{\{x,y\}\in X\times X\;\mid\;x\ne y\}
 $$ 
 is a given subset of unordered pairs of elements of $X$, and $m:X\to\Bbb R^+$. The elements of the set $X$ are called {\em vertices}, the elements of $E$ are called {\em edges}, and $m$ is the {\em measure}.
\end{definition}

For any $x,y\in X$ we write $x\sim y$ if $\{x,y\}\in E$. A graph is called {\em locally finite} if the set
$
\{y\;\mid\;y\sim x\}
$ is finite for all $x\in X$.  A {\em path} between  vertices $x,y\in X$ is a finite sequence $(x_0,\ldots,x_{n}), n\in \mathbb N\cup\{0\}$,  of vertices such that
$$
x=x_0\sim x_1\sim x_2\sim \dots \sim x_n=y.
$$

A set $ K\subseteq X $ is called {\em connected}, if there exists a path between any two vertices of it, and so we call  the graph \emph{connected} if $ X $ is connected.


\begin{assumption}\label{a:graph}  In this paper we assume that all graphs $ (X,E,m) $ are locally finite and connected.
\end{assumption}


Let us denote
$$
\Bbb H_r:=\{\cfreq \in \Bbb C\;\mid\;\Re \cfreq>0\}.
$$

\begin{definition}\label{def::complexweightedgraph}
A {\em complex-weighted graph} is a graph $(X,E,m)$, whose each edge is equipped with the weight $b:E\to \Bbb H_r$, satisfying the following property: there exists $c\ge 0$ such that
\begin{equation}\label{bsect}
\left|\Im b(x,y)\right|\le c \cdot \Re b(x,y).
\end{equation}
We call this property {\em sectoriality} and $c$ is a {\em sectoriality constant}.
\end{definition}
We assume $b(x,y)\equiv 0$ if there is no edge between $x$ and $y$. Hence, any complex-weighted graph is uniquely determined by the triple $(X,b,m)$, and we will refer to this triple as the complex-weighted graph.

\begin{assumption}
In this paper the complex weightsof a graph always satisfy a sectoriality \eqref{bsect}, i.e. all the weights of the graph belong to a sector of the complex plane.
\end{assumption}

The sectoriality  leads to many useful features of the graph and its energy form, the most basic of which are discussed in this section below. Moreover, the corresponding Dirichlet Laplacian is $m$-sectorial and generates a contractive holomorphic $C_0$-semigroup, see Section \ref{section::DirichletLaplacian}.



Let $X$ be an at most countable set and $(X,b,m)$ be a complex-weighted graph. We denote by $C(X)$ the set of all complex-valued functions on $X$.


The {\em formal Laplacian} is acting on  any function $f\in C(X)$ via
$$
\L_{b,m} f(x):=\dfrac{1}{m(x)}\sum_{y\in X} (f(x)-f(y)) b(x,y)
$$
We will omit subscripts $b$ and(or) $m$ when they are  clear from the context.



\subsection{Energy form and formal sesquilinear form }


 Let $\ell^2(X,m)$ be the set of functions
$$\{f\in C(X)\;\mid\;\sum_{x\in X}|f(x)|^2m(x)<\infty\big\},
$$
with a scalar product
$$
(f\mid g)=\sum_{x\in X} f(x)\overline{g(x)}m(x).
$$
Obviously, $\ell^2(X,m)$  is a Hilbert space.

Further, let us define a subspace $\mathcal D(X)\subset C(X)$ by

\begin{align*}
\mathcal D(X)&=\Big\{f\in C(X)\;\mid\;\sum_{x,y\in X}\abs{f(x)-f(y)}^2 |b(x,y)|<\infty\Big\}\\
&=\Big\{f\in C(X)\;\mid\; \sum_{x,y\in X}\abs{f(x)-f(y)}^2 b(x,y)\substack{ \mbox{ converges for some (all) }\\\mbox{ order of summations}}\Big\}\\
&=\Big\{f\in C(X)\;\mid\;\sum_{x,y \in X}\abs{f(x)-f(y)}^2 \Re b(x,y) <\infty\Big\},
\end{align*}
where the equalities follow from the sectoriality \eqref{bsect}.





The {\em (formal) energy} of a function is defined for any $f\in \mathcal D(X)$ as 
\begin{align*}
\mathcal Q_b(f)&=\dfrac12 \sum_{x,y\in X}\abs{f(x)-f(y)}^2b(x,y).
\end{align*}
We will omit the subscript $b$, if it is clear from the context.

%Note that $f\in\D(X)$ if and only if $\overline f\in \D(X)$.


Note that 
\begin{equation*}%\label{eq::Reqf}
\Re \mathcal Q(f)=\dfrac12 \sum_{x,y\in X}\abs{f(x)-f(y)}^2\Re b(x,y)\ge 0,
\end{equation*}
and
$$
\Im \mathcal Q(f)=\dfrac12 \sum_{x,y\in X}\abs{f(x)-f(y)}^2\Im b(x,y).
$$

Further,  the following lemma follows immediately from the sectoriality~\eqref{bsect}.
\begin{lemma}\label{lem::formalQmapstosector}
Let $(X,b,m)$ be a complex-weighted graph. The  form $\mathcal Q=\mathcal Q_b$ with the domain $\D(X)$ is sectorial, i.e. it satisfies
$$
|\Im \mathcal Q(f)|\le  c \cdot\Re \mathcal Q(f),
$$
for some $c\ge0 $ for any $f\in \D(X)$.
\end{lemma}
The next several lemmas are almost immediate corollaries of the definition of the form $\Q$ and are presented here for further references.

\begin{lemma}\label{lem::complexf}
Let $(X,b,m)$ be a complex-weighted graph. For any $f\in C(X)$ the following are equivalent:
\begin{itemize}
\item
$f\in \D(X)$,
 \item
 the both  functions $\Im f, \Re f\in \D(X)$.

 \end{itemize} 
 Moreover, if  any of the equivalent conditions is satisfied, the following holds:
\begin{equation}\label{eq::QfQRefQImf}
\mathcal Q(f)=\mathcal Q(\Re f)+\mathcal Q(\Im f).
\end{equation}
\end{lemma}
\begin{proof}
The statement follows from the sectoriality \eqref{bsect}, the fact that
$$
|f(x)-f(y)|^2=|\Re f(x)-\Re f(y)|^2+|\Im f(x)-\Im f(y)|^2.
$$
and the definition of $\mathcal Q$.
\end{proof}



The corresponding (formal) sesquilinear form is defined for any $f,g\in \D(X)$ by
\begin{align*}
\mathcal Q(f,g)&=\dfrac12 \sum_{x,y\in X}{(f(x)-f(y))}\overline{(g(x)-g(y))}b(x,y)
\end{align*}
The following inequality, which is a consequence of the to Cauchy-Schwarz inequality, is known for sectorial forms:

\begin{equation}\label{eq::CSforsect}
|\mathcal Q(f,g)|\le (1+c) \left(\Re \mathcal Q(f)\right)^\frac{1}{2} \left(\Re \mathcal Q(g)\right)^\frac{1}{2}.
\end{equation}



\bigskip
 Let us denote by $C_c(X)$ the subset of $C(X)$ of all complex-valued functions on $X$ with finite support.


\begin{lemma}
Let $(X,b,m)$ be a complex-weighted graph. Then $C_c(X)$ is dense in $\ell^2(X,m)$.
\end{lemma}

\begin{proof}
Let $f\in \ell^2(X,m)$. Let $(K_n)_{n\in \Bbb N}$ be a finite exhaustion of $(X, b, m)$, i.e. $K_n\subset X$ is finite and connected,  $K_{n}\subset K_{n+1}$ for any $n\in \Bbb N$ and $\cup_{n\in \Bbb N} K_n=X$. Let $f_n:=f\mid_{K_n}$, $n\in \Bbb N$. Then
$$
\|f-f_n\|^2=\sum_{x\in X\setminus K_n}|f(x)|^2m(x)\to 0,
$$
as $n\to\infty$.
\end{proof}



\begin{theorem}[Green's formula]
Let $(X,b,m)$ be a complex-weighted graph. Then 
\begin{enumerate}[leftmargin=*]
\item[{\em (a)}]
For all $f\in~\D(X)$ and $\phi \in C_c(X)$ the following holds
\begin{align*}
&\sum_{x\in X}\L f(x) \phi(x) m(x)=\sum_{x\in X}\L \phi(x) f(x)m(x)\\
&=\dfrac{1}{2}\sum_{x,y\in X}b(x,y)(\phi(x)-\phi(y))(f(x)-f(y)).
\end{align*}

\item[{\em (b)}]
For all $f\in \D\cap \ell^2(X,m)$ and $\phi \in C_c(X)$ the following holds
\begin{equation*}
\mathcal Q(\phi, f)=(\L \phi\mid f).
\end{equation*}
If, in addition, $\L f\in  \ell^2(X,m)$, then
\begin{equation*}
\mathcal Q(f,\phi)=( \L f\mid \phi).
\end{equation*}
\end{enumerate}
\end{theorem}

\begin{proof}
Due to the locally finitness of graph all the sums are finite. The equalities follow by direct computations, see e.g \cite[Lemma 7]{Muranova1}. 
\end{proof}




\subsection{Dirichlet Laplacian and corresponding form on infinite graphs}\label{section::DirichletLaplacian}


\begin{thebibliography}{99}
\bibitem[Muranova1]{Muranova1} Anna Muranova. \newblock On the notion of effective impedance. \newblock {\em Operator and Matrices}, 14(3):723--741 (2020). %
\end{thebibliography}
\end{document}
